% Nouveau EGA V (cote 25): archive cote 25, site page 14 (PDF page 15).
% typescript. Draft compared: none.
% AI-transcribed by Claude Opus 5.5 (claude-opus-5-5, adaptive thinking, effort high), Message Batches API, batch ega25b (msgbatch_01A7n5YuXubUPNK1zsepi8RG), 2026-11-13; not yet checked by eye.
\scanpage{25/14}
\mspage{\hand{11}}

ceci, soit Z l'ensemble des points en lesquels X n'est pas lisse sur
k, son image inverse \inserted{$Z_{y}$ dans $Y_{y}$} est donc (en vertu de 2.12.)   l'ensemble des
points de $Y_{\hand{y}}$ en lesquels $Y_{\hand{y}}$ n'estb pas lisse sur k($\hand{y}$). Or  la codimen-%
sion de  Z dans X en x  est égale à celle de  $Z_{\hand{y}}$ dans $Y_{\hand{y}}$  en y par
raison de platitude (réf par.6). Donc l'une est \hand{$>$}k  sss l'autre l'est,
ce qui achève la démonstration.
\uline{Corollaire} 2.15.   Soit Z l'ensemble des points de X en lesquels
X n'est pas lisse sur k (resp. n'est pas  géom $R_k$ sur k, resp. n'est
\margin{\unsure{//}} pas séparable sur \unsure{R|k} , resp. n'est pas \inserted{géom} normal sur k). Pour que $Y_{\hand{y}}$ soit
lisse (resp.  géom $R_k$, resp séparable, resp géom normal) sur k($\hand{y}$ ), il
f et s  que Z soit fini.

\margin{$|$} Du point da vue rédaction, 2.14. et 2.15.  contiennent 2.1\unsure{2|3}. et
2.13. (qu'on pourrait donc se dispenser d'énoncer séparément), d'autre
part le corollaire est plus important pratiquement que le théorème,
qu'on pour\unsure{r}ait donner en proposition ou lemme préliminaire, pour glori-%
fier plutôt le corollaire.

Nous pouvons  moduler \xout{\unsure{xxxxparxi}}  le cas de la propriété (iii):
\uline{Corollaire} 2.16.  (Supposant toutjours f une immersion \inserted{et de plus F cohérent}), sous les condi-%
\margin{\unsure{//}} tions de 2.7. , pour que $Z_i$ soit non immergé, il f et s que $Z_i$  le soit.
Lorsqu'il en est ainsi, la mu\hand{l}tiplicité radicielle de F en $Z_i$ \xout{(ref)}\deleted{rel} \inserted{rel.}
à k (réf) est égale à celle de $G_{\hand{y}}$  en $Z_i$  rel à k($\hand{y}$).

La première assertion est contenu dans la dernière assertion de
2.7. Pour la seconde, comme  $Y_{\hand{y}} \longrightarrow X$ est plat donc le foncteur F \hand{$\rightsquigarrow$} $G_{\hand{y}}$
exact,  on est ramené, par restriction à un voisinage du pt gén de $Z_i$
et dévissage, au cas où F = $\underline{O}_{Z_i}$ \inserted[typed]{et on peut supposer $Z_i = X$.} . D'ailleurs on peut commencer par
se réduire au cas k alg clos, \inserted[typed]{alors X est séparable sur k.} Alors \xout{\unsure{$\underline{O}$xxxxx$\underline{O}_{Z_i}$xxxxx}} l'assertion  est
contenue dans 2.15. (iii).   On en conclut comme d'hab itude:

\margin{14}
